\scanpage{47}
\margin{La construction de l'éclaté marche dans le contexte des variétés $C^\infty$ et \unsure{des variétés} d'ailleurs analytiques réelles.}
Pour ceci, considérons l'éclaté réel de $\Sigma$ en $J$, soit $\Sigma'$, dans lequel l'image inverse de $i\in J$ s'identifie à: \deleted{$\unsure{P(}$ l'espace} la droite projective réelle
\[
(11)\qquad\Sigma'_{R_i}\simeq P(\underbrace{T_{\Sigma,R_i}})
\]
\margin{\unsure{esp.|espace} tangent à $\Sigma$ en $R_i$, identifié à: $T_{\Sigma,R_i}^*/\mathbb R^*$}
qui est un cercle plongé dans la surface compacte $\Sigma'$ (qu'on peut considérer comme \unsure{obtenu|obtenue} en enlevant de la sphère trois rondelles autour des $R_i$, et y recollant des rubans de Möbius --- on obtient une surface non \uline{orientable} dont le \deleted{genre} $\chi_!$ est égal à
\[
\begin{aligned}
&\chi_!(\Sigma)\deleted{\unsure{-3}}\inserted{-3\chi_!(\text{rondelle ouverte})}+3(\underbrace{\chi_!(\text{ruban Möbius})}_{0}\\
&\deleted{\unsure{\text{à}}}=2-3=-1=1-g,
\end{aligned}
\]
donc le genre est $2$).

La découpe de cette surface suivant $\bigcup\Sigma'_{R_i}$ est une surface à bord compacte, \deleted{dont} \inserted{en fait} soit $\widetilde\Sigma$, avec
\margin{$\partial\widetilde\Sigma$ un rev. double de $\coprod_i\Sigma'_{R_i}$}
\[
(12)\qquad\partial\widetilde\Sigma\simeq\coprod_{i\in I}\underbrace{TU_{\Sigma,R_i}}
\]
\margin{cercle tangent à $\Sigma$ en $R_i$, identifiable à $T_{\Sigma,R_i}^*/\mathbb R^{*+}$ noté aussi $U_i$}
\margin{$\widetilde\Sigma$ est une sphère à bord à trois trous. Notons que $\widetilde\Sigma^\circ$ a une structure conforme, mais que celle-ci ne se prolonge pas à $\widetilde\Sigma$.}
\margin{46}
